\documentclass[11pt,a4paper]{article}

\usepackage[margin=2.2cm]{geometry}
\usepackage{amsmath,amssymb,amsthm,mathtools}
\usepackage[T1]{fontenc}
\usepackage[dvipsnames]{xcolor}
\usepackage{hyperref}
\usepackage{enumitem}
\usepackage[most]{tcolorbox}
\usepackage{tikz}

\hypersetup{
  colorlinks=true,
  linkcolor=blue!55!black,
  urlcolor=blue!55!black
}

\setlist[enumerate]{itemsep=0.35em, topsep=0.45em}

\newtheorem{theorem}{Theorem}[section]
\newtheorem{definition}[theorem]{Definition}
\newtheorem{example}[theorem]{Example}
\newtheorem{remark}[theorem]{Remark}

\tcbset{
  commonbox/.style={
    enhanced,
    breakable,
    boxrule=0.8pt,
    arc=2mm,
    left=1.4mm,
    right=1.4mm,
    top=1mm,
    bottom=1mm,
    colback=white
  }
}

\newtcolorbox{problemblock}{
  commonbox,
  colframe=BrickRed!65!black,
  colback=BrickRed!4,
  title=Problem,
  fonttitle=\bfseries
}

\newtcolorbox{intuitionblock}{
  commonbox,
  colframe=RoyalBlue!70!black,
  colback=RoyalBlue!4,
  title=Intuition,
  fonttitle=\bfseries
}

\newtcolorbox{solutionblock}{
  commonbox,
  colframe=ForestGreen!60!black,
  colback=ForestGreen!4,
  title=Solution,
  fonttitle=\bfseries
}

\newtcolorbox{keypointblock}{
  commonbox,
  colframe=black!65,
  colback=black!3,
  title=Key Point,
  fonttitle=\bfseries
}

\newtcolorbox{mathblock}[1]{
  commonbox,
  colframe=Purple!60!black,
  colback=Purple!4,
  title=#1,
  fonttitle=\bfseries
}

\newcommand{\vect}[1]{\boldsymbol{#1}}

\title{[Topic Title]\\\large [Optional Subtitle]}
\author{}
\date{\today}

\begin{document}

\maketitle
\tableofcontents
\newpage

\section{What This Lesson Is Trying to Explain}

\begin{problemblock}
[State the real-world or learning problem first.]
\end{problemblock}

[Open with short, plain sentences. If the user already gave a good intuitive paragraph,
reuse and polish it here instead of replacing it.]

\begin{intuitionblock}
[Give the first mental picture before introducing formal notation.]
\end{intuitionblock}

\begin{keypointblock}
[Write the one-sentence takeaway the beginner should keep in mind.]
\end{keypointblock}

\section{Basic Definition}

\begin{mathblock}{Definition}
\begin{definition}
[Insert the core definition.]
\end{definition}
\end{mathblock}

[Explain each new word immediately in ordinary language.]

\section{The First Picture or Smallest Example}

[Add a tiny numerical example, coordinate picture, or one TikZ figure if the idea is geometric.]

\section{A Basic Fact or Theorem}

\begin{mathblock}{Theorem}
\begin{theorem}
[Insert one basic fact worth remembering.]
\end{theorem}
\end{mathblock}

\begin{mathblock}{Proof}
\begin{proof}
[Keep the proof short. For beginners, prefer a geometric or arithmetic proof over abstraction.]
\end{proof}
\end{mathblock}

\begin{mathblock}{Example}
\begin{example}
[Work one example completely.]
\end{example}
\end{mathblock}

\section{Connection to Physics or Code}

\begin{problemblock}
[Ask where this idea actually appears.]
\end{problemblock}

\begin{solutionblock}
[Connect the idea to physics, code, or the user’s actual context.]
\end{solutionblock}

\section{What to Remember}

\begin{enumerate}[label=\arabic*.]
  \item [Key point 1]
  \item [Key point 2]
  \item [Key point 3]
\end{enumerate}

\begin{keypointblock}
[End with one sentence simple enough to remember.]
\end{keypointblock}

\end{document}
